\documentclass[11pt,letterpaper]{article}

% Match website font: PT Serif
\usepackage[T1]{fontenc}
\usepackage{PTSerif}
\usepackage[letterspace=120]{microtype}

% Full-width layout: section headers sit above their content
\usepackage[
  letterpaper,
  left=0.8in,
  right=0.8in,
  top=0.95in,
  bottom=1.0in,
]{geometry}

\usepackage{xcolor}
\usepackage{hyperref}
\usepackage{enumitem}
\usepackage{fancyhdr}
\usepackage{ragged2e}
% Keep a header (or an entry title) from being stranded at the foot of a page.
% Hand-rolled instead of \needspace, whose \vfil ejects half-empty pages here.
\newdimen\cvroom
\newcommand{\cvneed}[1]{%
  \par
  \ifdim\pagegoal=\maxdimen\else
    \cvroom=\dimexpr\pagegoal-\pagetotal\relax
    \ifdim\cvroom<#1\relax\newpage\fi
  \fi}

\definecolor{labelgray}{gray}{0.42}
\definecolor{rulegray}{gray}{0.72}
\definecolor{linkblue}{HTML}{1A4B8C}
\hypersetup{colorlinks=true, urlcolor=linkblue, linkcolor=linkblue}

\pagestyle{fancy}
\fancyhf{}
\renewcommand{\headrulewidth}{0pt}
\fancyfoot[R]{\small\itshape Curriculum Vitae, Minjin Chae, \thepage}

% ---- Section header: full-width, above the content ----
\newcommand{\cvsection}[1]{%
  \par\addvspace{1.15em}%
  \cvneed{4\baselineskip}%
  \noindent{\small\bfseries\textls{\MakeUppercase{#1}}}\par
  \vspace{0.25em}%
  {\color{rulegray}\hrule height 0.5pt}%
  \vspace{0.55em}%
  \nobreak\ignorespaces}

% ---- Subline: indented, italic, gray (awards attached to an entry) ----
% NOTE: \RaggedRight must come BEFORE \leftskip -- ragged2e resets \leftskip to 0.
% One paragraph with \interlinepenalty=10000, so a multi-line subline can never
% be split across a page; \nobreak keeps it with the entry it belongs to.
% Separate multiple sublines with \\ (not a blank line).
% ---- Status: inline tag at the end of an entry, not on a line of its own ----
\newcommand{\status}[1]{\hspace{0.5em}{\small\itshape\color{labelgray}#1}}

\newcommand{\subline}[1]{%
  \par\nobreak\vspace{0.15em}%
  \begingroup
  \small\itshape
  \RaggedRight
  \interlinepenalty=10000
  \setlength{\parindent}{0pt}%
  % \leavevmode BEFORE \color: a \color whatsit in vertical mode would make the
  % following \parskip glue a legal page break, detaching this from its entry.
  \leavevmode\color{labelgray}%
  #1\par
  \endgroup
  \vspace{0.3em}}

% ---- Sublines: indented block of plain lines (venues, courses) ----
\newenvironment{sublines}{%
  \par\nobreak
  \begingroup
  \small
  \RaggedRight
  \leftskip=1.6em
  \setlength{\parskip}{0.32em}\setlength{\parindent}{0pt}%
}{\par\endgroup\vspace{0.2em}}

% ---- Entry: a titled block (paper, presentation topic) ----
\newcommand{\entry}[1]{\par\addvspace{0.85em}\cvneed{3\baselineskip}\noindent #1\par\nobreak}

\setlist{leftmargin=1.5em, topsep=1pt, partopsep=0pt, itemsep=1pt, parsep=0pt}
\setlength{\parindent}{0pt}
\setlength{\parskip}{0.45em}
\setlength{\emergencystretch}{2em}
\raggedbottom
\widowpenalty=10000
\clubpenalty=10000

% True ragged right, no hyphenation (cleaner for short CV lines)
\setlength{\RaggedRightRightskip}{0pt plus 1fil}
\hyphenpenalty=10000
\exhyphenpenalty=10000

\begin{document}
\RaggedRight

{\huge Minjin Chae}\par
\smallskip
\href{mailto:minjinchae@g.harvard.edu}{minjinchae@g.harvard.edu} \enspace\textbar\enspace \url{https://chaeminjin.github.io/}\par
\medskip

\cvsection{Education}
\textbf{Harvard University}\\
Ph.D. in Sociology, 2027 (expected)\\
Dissertation: \emph{Flexible Work, Durable Inequality: Gender, Race, and the Remaking of Workplace Advantage}\\
Committee: David Pedulla (chair), Daniel Schneider, Alexandra Killewald (University of Michigan)\\
M.A. in Sociology, 2023

\smallskip
\textbf{Seoul National University}\\
M.A. in Sociology, 2019\\
B.A. in History Education and B.B.A. in Business, 2017

\cvsection{Research interests}
future of work, flexible work, social inequality, work and family, gender and race

\cvsection{Publications}
Yun, Youngho, Yeani Choi, and Minjin Chae. 2021. ``Development and Validity Testing of the Workplace Parentship Index (WPI): Assessment of Family-Friendly Practices.'' \emph{Journal of Occupational and Environmental Medicine} 63(12).

\cvsection{Working papers}
Chae, Minjin. ``When Does Gender Trump Flexibility? Workplace Flexibility and Domestic Task Assignment.'' Under review.
\subline{Finalist, Best Student Paper Award, Organization and Management Theory Division, Academy of Management, 2026}

\cvsection{Work in progress}
Chae, Minjin. ``Activating Inequality: Task Segregation and the Racial and Gender Divide in Remote Work.''\status{Draft available upon request}

Chae, Minjin. ``Worker Power and Temporal Rents: How Occupational Closure and Productivity Constraints Shape Schedule Control.''\status{Draft available upon request}

Cha, Youngjoo, Kaitlin Johnson, and Minjin Chae. ``It's the Hours, Not the Job: Temporal Flexibility and the Gender Wage Gap Within Occupations.''\status{Draft available upon request}

Chae, Minjin. ``When Do Firms Reduce or Reproduce Gender Inequality? Organizational Contexts of the Pandemic Work From Home and KiF\"oG Childcare Expansions.''\status{Data analysis}
\subline{Guest Researcher in residence, WZB Berlin Social Science Center, Spring 2026}

Chae, Minjin. ``Who Works From Home Within the Same Job? Racial Inequality in Remote Work and the Role of Collective Bargaining.''\status{Data analysis}

Chae, Minjin. ``Gender Essentialist Norms and Skeptical Attitudes Toward Marriage Among South Korean Young Adults.''\status{Draft available upon request}
\subline{International Conference on Sociology of Korea Student Paper Award, University of Pennsylvania Kim Center for Korean Studies, 2022\\
Best Graduate Student Paper Award, The Population Association of Korea, 2019}

\cvsection{Presentations}
\entry{\textit{Occupational Segregation and Remote Work Inequality by Race and Gender, 2017--2026}}
\begin{sublines}
American Sociological Association Annual Meeting. New York City, US. August 2026.

Research Committee 28 of the International Sociology Association. New York City, US. August 2026.
\end{sublines}

\entry{\textit{Can Workplace Flexibility Policies Promote Gender Equality at Home?}}
\begin{sublines}
Academy of Management Annual Meeting. Philadelphia, US. August 2026.

Work and Family Researcher Network Conference. Montreal, Canada. June 2026.

Research Committee 28 of the International Sociology Association. Seville, Spain. May 2026.

American Sociological Association Annual Meeting. Chicago, US. August 2025.

International Community, Work and Family Conference. Bielefeld, Germany. March 2025.

Korean Sociological Association Annual Meeting. Seoul, South Korea. December 2024.
\end{sublines}

\entry{\textit{How Worker Power and Time-related Productivity Constraints Shape Schedule Quality}}
\begin{sublines}
American Sociological Association Annual Meeting. Montreal, Canada. August 2024.

American Sociological Association Annual Meeting. Philadelphia, US. August 2023.

Korea Inequality Research Network Colloquium. Chuncheon, South Korea. July 2023.
\end{sublines}

\entry{\textit{The Role of Work Hours and Occupation in the Gender Wage Gap}}
\begin{sublines}
Work and Family Researcher Network Conference. Montreal, Canada. June 2024.
\end{sublines}

\entry{\textit{Gender Essentialist Norms and Skeptical Attitudes Toward Marriage Among South Korean Young Adults}}
\begin{sublines}
International Conference on Sociology of Korea. Philadelphia, US. October 2022.

Population Association of Korea Annual Meeting. Daejeon, South Korea. June 2019.
\end{sublines}

\cvsection{Honors \& awards}
Finalist, Best Student Paper Award, Organization and Management Theory Division, Academy of Management, 2026

Doctoral Consortium, Organization and Management Theory Division, Academy of Management, 2026

Teaching Certificate, Derek Bok Center for Teaching and Learning, Harvard University, 2025

International Conference on Sociology of Korea Student Paper Award, University of Pennsylvania Kim Center for Korean Studies, 2022

Best Graduate Student Paper Award, The Population Association of Korea, 2019

2nd Prize, Student Paper Competition, Institute for Peace and Unification Studies, Seoul National University, 2017

Grand Prize, Undergraduate Student Paper Competition, The Korean Society for Chinese Studies, 2014

\cvsection{Fellowships \& grants}
Dissertation Completion Fellowship, Harvard University Graduate School of Arts and Sciences, 2026--2027

Graduate Seed Grant, Harvard University Center for American Political Studies, 2026

Summer Research Grant, Harvard Griffin GSAS Student Council, 2026

Krupp Foundation Dissertation Research Fellowship, Minda de Gunzburg Center for European Studies, Harvard University, 2026

Travel Grant, Harvard Center for Population and Development Studies, 2024--2026

Travel Grant, Harvard University Department of Sociology, 2022--2026

Merit/Graduate Society Term-Time Research Fellowship, Harvard University Graduate School of Arts and Sciences, 2025

Jens Aubrey Westengard Fund, Harvard University Graduate School of Arts and Sciences, 2025

Student Scholarship, Seoul National University Alumni Association of New England, 2025

Dissertation Research Grant, Harvard Center for Population and Development Studies, 2024

Research Grant, Harvard University Institute for Quantitative Social Science, 2024

Travel Grant, Harvard University Institute for Quantitative Social Science, 2024, 2025

Graduate Seed Grant, Harvard University Center for American Political Studies, 2023

Summer Institute in Computational Social Science Research Grant, Korea Development Institute School of Public Policy and Management, 2022

Division of Social Science Fund, Harvard University Graduate School of Arts and Sciences, 2022

Community Researcher Grant, Seoul Metropolitan Government, 2019

Lecture and Research Scholarship, Seoul National University Department of Sociology, 2018

Naum Scholarship, Seoul National University Department of Sociology, 2017

SNU Undergraduate Research Program Grant, Seoul National University, 2016

Study Abroad Program Grant, Seoul National University Office of International Affairs, 2016

Merit-Based Scholarship, Seoul National University, 2013--2015

SNU Development Fund Scholarship, Seoul National University, 2013

\cvsection{Teaching}
\entry{\textbf{Teaching Fellow, Harvard University}}
\begin{sublines}
Work and Culture. Fall 2024.

Models of Social Science Research. Spring 2022, 2023.

Introduction to Sociology. Fall 2022, 2023.
\end{sublines}

\entry{\textbf{Undergraduate Thesis Advisor, Harvard University}}
\begin{sublines}
Jenna Hong

Mira Nalbandian
\end{sublines}

\cvsection{Service}
Committee on Higher Degrees, Department of Sociology, Harvard University (Member), 2025--2026

Department Colloquium Committee, Department of Sociology, Harvard University (Member), 2024--2025

Graduate Student Organization, Department of Sociology, Harvard University (Social Chair), 2022--2024

Memberships: American Sociological Association, Academy of Management, Work Family Researchers Network, Research Committee 28 of the International Sociology Association

\cvsection{Professional experience}
\entry{\textbf{Research Associate, Deokinwon Co., Ltd., Seoul, South Korea}, 2020}
\begin{sublines}
Developed and validated the Workplace Parentship Index, an instrument assessing employee wellbeing, health, and family-friendly workplace practices.

Fielded employee and HR manager surveys with a client firm and drafted the assessment report.
\end{sublines}

\cvsection{Technical skills}
Stata, R, LaTeX, Qualtrics

\cvsection{Languages}
English (Fluent), Korean (Native), German (Intermediate, B1), Mandarin (Beginner)

\cvsection{References}
{\small
\begin{tabular}[t]{@{}p{0.30\linewidth}p{0.30\linewidth}p{0.30\linewidth}@{}}
David Pedulla & Daniel Schneider & Alexandra Killewald \\
Department of Sociology & Department of Sociology & Department of Sociology \\
Harvard University & Harvard University & University of Michigan \\
\href{mailto:dpedulla@fas.harvard.edu}{dpedulla@fas.harvard.edu} & \href{mailto:dschneider@hks.harvard.edu}{dschneider@hks.harvard.edu} & \href{mailto:akillewa@umich.edu}{akillewa@umich.edu} \\
\end{tabular}}

\end{document}
